% ---------------------------------------------------
%
% MODERATION CARDS, DIN A5 landscape, cover plus one card per page
%
% ---------------------------------------------------
%
% MMK Heilbronn, 24 November 2026, 14:15-14:45
% Panel "Can Industrial AI Turn Germany's Weaknesses into Strengths?"
%
% Nine cards to hold on stage: two opening cards, three impulse speakers,
% discussion guest Eike Wibrow (no impulse), two thesis pairs, and the outlook.
% Everything on a card is meant to be read at arm's length under stage light:
% few words, large type, the spoken lines in English.
% The full version with sources and reserve questions is "moderation deck.md".
%
% Build from the repository root (the PDF lands in build/):
%   latexmk -pdf -outdir=build -auxdir=build_aux "01 MMK 2026 Panel/20261124_MMK_Moderationskarten_A5.tex"
% Export the cover separately for use as the outward-facing card/backing:
%   pdfseparate -f 1 -l 1 build/20261124_MMK_Moderationskarten_A5.pdf build/20261124_MMK_Titelseite_A5.pdf
% Export Wibrow's introduction separately:
%   pdfseparate -f 7 -l 7 build/20261124_MMK_Moderationskarten_A5.pdf build/20261124_MMK_Eike_Wibrow_A5.pdf
% Wuttke's own speaking cards are PDF pages 2, 3 and 10.
% Print at actual size on A5 landscape, or two pages per A4 portrait sheet.

\documentclass[12pt]{article}

\usepackage[a5paper,landscape,left=10mm,right=10mm,top=9mm,bottom=9mm]{geometry}
\usepackage{iftex}
\ifPDFTeX
  \usepackage[utf8]{inputenc}
  \usepackage[T1]{fontenc}
\fi
\usepackage{helvet}
\renewcommand{\familydefault}{\sfdefault}
\usepackage{xcolor}
\usepackage{graphicx}
\graphicspath{{figures/}{../figures/}}
\usepackage{tikz}
\usetikzlibrary{calc}
\usepackage{enumitem}
\usepackage{microtype}

\pagestyle{empty}
\setlength{\parindent}{0pt}
\setlength{\parskip}{0pt}
\raggedright
\frenchspacing
\hyphenpenalty=10000 \exhyphenpenalty=10000   % cue cards: never break a word

% TUM colours, as in the slide deck: A is blue, B is orange.
\definecolor{primaryBlue}{RGB}{0,101,189}
\definecolor{tertiaryOrange}{RGB}{227,114,34}
\definecolor{secondaryGray}{RGB}{102,102,102}

% ---------------------------------------------------
% Card elements
% ---------------------------------------------------

% Cover/backing: white margins allow printing without a borderless printer.
\newcommand{\coverpage}{%
  \noindent\begin{minipage}[t]{0.65\linewidth}
    \vspace{0pt}
    {\color{primaryBlue}\fontsize{12}{15}\selectfont\bfseries
      TUM School of Management\par}
    \vspace{1mm}
    {\color{secondaryGray}\fontsize{11}{14}\selectfont TUM Campus Heilbronn\par}
  \end{minipage}\hfill
  \begin{minipage}[t]{40mm}
    \vspace{0pt}\includegraphics[width=\linewidth]{TUM_Logo_blau_transparent.png}
  \end{minipage}\par
  \vfill
  {\color{primaryBlue}\fontsize{11}{14}\selectfont\bfseries MMK HEILBRONN 2026\par}
  \vspace{4mm}
  {\color{primaryBlue}\fontsize{27}{32}\selectfont\bfseries
    Can Industrial AI Turn\\
    Germany's Weaknesses\\
    into Strengths?\par}
  \vspace{5mm}
  {\color{primaryBlue}\rule{28mm}{1.2pt}\par}
  \vfill
  {\fontsize{12}{15}\selectfont\bfseries Prof. Dr. David Wuttke\par}
  \vspace{1.5mm}
  {\color{secondaryGray}\fontsize{10}{13}\selectfont
    Moderation\hfill 24 November 2026\par}
  \newpage}

% Two fixed columns keep each card together on one landscape page.
\newcommand{\cardcolumns}[2]{%
  \noindent\begin{minipage}[t]{0.45\linewidth}\vspace{0pt}\raggedright
    #1
  \end{minipage}\hfill
  \begin{minipage}[t]{0.51\linewidth}\vspace{0pt}\raggedright
    #2
  \end{minipage}\par}

% Header band: #1 kicker (left), #2 time (right), #3 name or title, #4 role or subtitle
\newcommand{\cardhead}[4]{%
\noindent
\begin{tikzpicture}
\node[fill=primaryBlue, text=white, inner xsep=4mm, inner ysep=2.4mm, outer sep=0pt,
      text width=\dimexpr\linewidth-8mm\relax, align=left] {%
  {\fontsize{9}{11}\selectfont\bfseries #1\hfill #2\par}\vspace{1mm}
  {\fontsize{21}{23}\selectfont\bfseries #3\par}\vspace{0.6mm}
  {\fontsize{11}{13}\selectfont #4\par}};
\end{tikzpicture}\par}

% Small blue label above each block
\newcommand{\lbl}[1]{\par\vspace{3mm}%
  {\fontsize{8.5}{10}\selectfont\bfseries\color{primaryBlue}\MakeUppercase{#1}\par}\vspace{1mm}}

% Two or three facts about the speaker
\newenvironment{facts}{%
  \begin{itemize}[leftmargin=4.5mm, labelsep=2mm, itemsep=1.4mm, topsep=0mm, parsep=0mm,
                  partopsep=0mm, label={\color{primaryBlue}\textbullet}]%
  \fontsize{12}{15}\selectfont}{\end{itemize}}

% Larger keywords for the moderator's own opening and outlook.
\newenvironment{cues}{%
  \begin{itemize}[leftmargin=4.5mm, labelsep=2mm, itemsep=2mm, topsep=0mm, parsep=0mm,
                  partopsep=0mm, label={\color{primaryBlue}\textbullet}]%
  \fontsize{13.5}{17}\selectfont}{\end{itemize}}

% Question to the room: #1 question, #2 what to ask hands for
\newcommand{\ask}[2]{%
\par\noindent
\begin{tikzpicture}
\node[fill=primaryBlue!8, inner xsep=4mm, inner ysep=3mm, outer sep=0pt,
      text width=\dimexpr\linewidth-8mm\relax, align=flush left] (q) {%
  {\fontsize{15}{18.5}\selectfont #1\par}\vspace{2.6mm}
  {\fontsize{11.5}{14}\selectfont\bfseries\color{primaryBlue}#2\par}};
\draw[primaryBlue, line width=1.4mm] ([xshift=0.7mm]q.north west) -- ([xshift=0.7mm]q.south west);
\end{tikzpicture}\par}

% The line that hands over to the speaker
\newcommand{\handover}[1]{\par{\fontsize{15}{18.5}\selectfont\bfseries #1\par}}

% Plain note, e.g. what to listen for
\newcommand{\note}[1]{\par{\fontsize{11.5}{14.5}\selectfont #1\par}}

% Footer and page break
\newcommand{\cardfoot}[1]{\par\vspace{2mm}\vfill
  \noindent\begin{minipage}{\linewidth}
    \color{secondaryGray}\fontsize{9}{11}\selectfont
    \rule{\linewidth}{0.4pt}\par\vspace{1.2mm}
    #1\par
  \end{minipage}\par\newpage}

% Thesis as on the slide: #1 colour, #2 letter, #3 text
\newcommand{\thesis}[3]{%
\par\noindent
\begin{tikzpicture}
\node[fill=#1!10, inner xsep=3.2mm, inner ysep=2.2mm, outer sep=0pt, anchor=north west,
      text width=\dimexpr\linewidth-9mm-6.4mm\relax, align=flush left,
      font=\fontsize{11.5}{14.5}\selectfont] (t) at (9mm,0) {#3};
\fill[#1] (0,0) rectangle (t.south west);
\node[text=white, inner sep=0pt, font=\bfseries\fontsize{17}{17}\selectfont]
      at ($(0,0)!0.5!(t.south west)$) {#2};
\end{tikzpicture}\par\vspace{1mm}}

% Question to one panelist on a pair card: #1 name, #2 question
\newcommand{\q}[2]{\par\vspace{2mm}%
  {\fontsize{11.5}{14.5}\selectfont{\bfseries\color{primaryBlue}#1}\enspace #2\par}}

\newcommand{\then}{\,{\fontsize{12}{12}\selectfont$\rightarrow$}\,}

\begin{document}

\coverpage

% ---------------------------------------------------
% CARD 1: Welcome and focus, based on Wuttke's notes of 23 November 2026
% ---------------------------------------------------
\cardhead{CARD 1 OF 9 \textbar{} OPENING}{14:16 \textbar{} 1 MIN}{Welcome to Heilbronn}{Bildungscampus, skilled labour and physical AI}

\cardcolumns{
\lbl{Welcome}
\handover{Wonderful to see you here in Heilbronn!}

\lbl{Welcome}
\begin{cues}
\item Many strengths and weaknesses - one particularly relevant: skilled labor
\item Some say strength, some say weakness. 
\item How will physical AI change this tension? Let's figure this our and learn from our panel!
\end{cues}
}

\lbl{The question}
\handover{Skilled labour: a strength, a weakness --- and what can physical AI change?}

\cardfoot{Next: AR Study}

% ---------------------------------------------------
% CARD 2: Research and transition
% Study findings are supplied by Wuttke in his notes of 23 November 2026.
% AR is the research example informing questions about physical AI;
% no effect sizes or findings about all forms of physical AI are implied.
% ---------------------------------------------------
\cardhead{CARD 2 OF 9 \textbar{} RESEARCH}{14:17 \textbar{} 3 MIN}{Assistance and learning}{AR as implementation of physical AI}

\cardcolumns{
\lbl{Study 1: AR guidance}
\begin{cues}
\item How does physical AI through AR affect processes? What happens with your skilled labor when you introduce AR?
\item Field experiment
\item If done well, AR makes your skilled labor more efficient. But what about learning and optmization?
\end{cues}

\lbl{Transition, by 14:20}
\handover{Three short impulses, then two pairs of opposing theses.}
\vspace{1mm}
\note{Ahr\then Engel\then Grunwald. Wibrow joins the discussion.}
}

\cardfoot{Next: introduce Helmut Ahr, 14:20 \textbar{} AR findings as a starting point for the physical AI discussion}

% ---------------------------------------------------
% CARD 3: Helmut Ahr
% ---------------------------------------------------
\cardhead{CARD 3 OF 9 \textbar{} IMPULSE 1}{14:20 \textbar{} 4 MIN}{Helmut Ahr}{CEO, Horváth}

\cardcolumns{
\lbl{Who}
\begin{facts}
\item 26 years at Horváth, CEO since 2022
\item Leads Technology \& AI consulting himself
\end{facts}

\lbl{Listen for}
\note{1 humanoid for 2.5 workers \textbar{} 10 euros per hour \textbar{} 50\,\% of tasks by 2040.\\[1mm]
Six slides in four minutes: watch the timer.}
}{
\lbl{Ask the room}
\ask{Humanoid robots are entering the shop floor. Would you feel you should be among the first, or rather wait until we have more insights from peers?}
    {Hands: among the first? \qquad Rather wait?}

\lbl{Hand over}
\handover{``You should be fast,'' says Helmut Ahr.}
}

\cardfoot{Next: Sarah Engel, 14:24}

% ---------------------------------------------------
% CARD 4: Sarah Engel
% ---------------------------------------------------
\cardhead{CARD 4 OF 9 \textbar{} IMPULSE 2}{14:24 \textbar{} 4 MIN}{Sarah Engel}{Head of AI, TRUMPF}

\cardcolumns{
\lbl{Who}
\begin{facts}
\item Leadership roles at IBM and Capgemini: took companies from their first AI strategy to the enterprise-wide rollout
\item Studied cognitive science and neuroinformatics in Tübingen and Ulm
\item Close to home: TRUMPF joined IPAI here in Heilbronn this July
\end{facts}

\lbl{Listen for}
\note{The Cutting Assistant: what happened to the operator's skills, and how did TRUMPF measure it?}
}{
\lbl{Ask the room}
\ask{An AI assistant suggests the settings for your most expensive machine. Do you press start, or do you ask your most experienced operator first?}
    {Hands: press start? \qquad Ask the expert first?}

\lbl{Hand over}
\handover{AI that delivers real value, and that people trust: that is Sarah Engel's job at TRUMPF.}
}

\cardfoot{Next: Michael Grunwald, 14:28}

% ---------------------------------------------------
% CARD 5: Michael Grunwald
% ---------------------------------------------------
\cardhead{CARD 5 OF 9 \textbar{} IMPULSE 3}{14:28 \textbar{} 4 MIN}{Michael Grunwald}{Head of Architecture \& Cross-Technologies, CARIAD}

\cardcolumns{
\lbl{Who}
\begin{facts}
\item Software and system architecture at CARIAD, the Volkswagen Group's software company
\item Also vehicle platforms, over-the-air technologies and AI-driven engineering
\item His organisation includes CARIAD's R\&D AI Factory: AI agents in automotive software development
\end{facts}

\lbl{Listen for}
\note{Who does he still need on his teams? Who becomes the next architect?}
}{
\lbl{Ask the room}
\ask{Software in five years: is it built by large teams of engineers, or by a few top people steering a fleet of AI agents?}
    {Hands: large teams? \qquad A few, with agents?}

\lbl{Hand over}
\handover{Michael Grunwald's AI Factory at CARIAD is finding out right now.}
}

\cardfoot{Next: introduce Eike Wibrow, then discussion, 14:32}

% ---------------------------------------------------
% CARD 6: Eike Wibrow, discussion only (no separate impulse)
% Sources, checked 23 November 2026: Wibrow (KION)/Eike Wibrow - Profil MMK.md
% KION, "Strong Partnerships, Smarter Warehouses", 21 October 2026;
% Wibrow's own CES 2025 post; KION's start-up article, 12 June 2025.
% GXO deployment metric: KION press release, 21 October 2026.
% This is commissioning time, not a workforce reduction.
% ---------------------------------------------------
\cardhead{CARD 6 OF 9 \textbar{} DISCUSSION GUEST}{14:32 \textbar{} SHORT INTRO}{Eike Wibrow}{VP Global Technology Strategy, KION Group}

\cardcolumns{
\lbl{Who}
\begin{facts}
\item Helped launch KION's partnership with NVIDIA and Accenture
\item Focus: autonomous warehouses, digital twins and scaling automation
\item Builds on technology partnerships and cooperation with start-ups
\end{facts}

\lbl{Concrete example}
\note{GXO, \'Epinoy: autonomous vehicles work alongside warehouse staff.\\[1mm]
KION reports over 60\,\% less on-site deployment time through simulation.}

\lbl{Listen for}
\note{Which tasks disappear? Who handles exceptions? How do beginners learn?}
}{
\lbl{Introduce, about 20 seconds}
\handover{Eike Wibrow leads global technology strategy at KION. He brings autonomous warehouses and physical AI into our discussion about people and skills.}

\lbl{Bridge to pair 1}
\note{From the shop floor to the warehouse: are we replacing skilled people, or preserving their experience?}

\lbl{If time, during discussion}
\q{Wibrow}{At GXO, what changed for the people: their headcount, their tasks, or their skills?}
}

\cardfoot{Next: pair 1 \textbar{} Introduction is part of the 14:32--14:38 discussion slot}

% ---------------------------------------------------
% CARD 7: Pair 1
% ---------------------------------------------------
\cardhead{CARD 7 OF 9 \textbar{} PAIR 1 \textbar{} SLIDE 1/2}{14:32 \textbar{} 6 MIN}{The boomers}{What is industrial AI for?}

\cardcolumns{
\vspace{4mm}
\thesis{primaryBlue}{A}{If industrial AI helps, it is because we can finally afford \textbf{fewer expensive skilled workers}.}
\thesis{tertiaryOrange}{B}{If industrial AI helps, it is because it catches the \textbf{boomers' knowledge} before it retires.}

\lbl{How}
\note{Hands in the room: A or B? Then \textbf{Ahr\then Engel\then Grunwald\then Wibrow}.\\
A or B, one reason, 45 seconds. No ``both are right''.}

\lbl{If time}
\note{To Ahr: the boomers are gone by 2037, your robots reach 50 percent after 2040. Who runs the lines in between?}
}{
\lbl{One question each}
\q{Ahr}{One humanoid for 2.5 workers. But your robots learn from today's skilled workers. So: replace them, or learn from them before they retire?}
\q{Engel}{The Cutting Assistant needs no cutting experience. Is that A or B: fewer experienced operators, or keeping what they know?}
\q{Grunwald}{30 percent leaner, 1,000 AI experts in Berlin. Fewer expensive engineers, or knowledge brought back in-house?}
\q{Wibrow}{Autonomous trucks: fewer drivers, or preserving their know-how? Which experienced judgement is hardest to capture?}
}

\cardfoot{Next: pair 2, 14:38}

% ---------------------------------------------------
% CARD 8: Pair 2
% ---------------------------------------------------
\cardhead{CARD 8 OF 9 \textbar{} PAIR 2 \textbar{} SLIDE 2/2}{14:38 \textbar{} 6 MIN}{Gen Z}{What does industrial AI do to them?}

\cardcolumns{
\vspace{4mm}
\thesis{primaryBlue}{A}{For Gen Z, industrial AI is a \textbf{shortcut} to expertise. They start where the boomers ended.}
\thesis{tertiaryOrange}{B}{For Gen Z, industrial AI is a \textbf{substitute} for expertise. They never have to learn what the machine already knows and will never become as much experts as boomers are.}

\lbl{How}
\note{Hands in the room: A or B? Then \textbf{Engel\then Grunwald\then Ahr\then Wibrow}.\\
45 seconds each.}
}{
\lbl{One question each}
\q{Engel}{A beginner cuts good parts on day one. Two years later: an expert, or a beginner with a very good assistant?}
\q{Grunwald}{If AI writes the simple code, who becomes your next architect?}
\q{Ahr}{Your roadmap reads like an apprenticeship: the robot takes the simple tasks first. Where does someone who starts in 2030 learn to handle a disruption?}
\q{Wibrow}{When a truck handles routine moves, how does a beginner learn to deal with exceptions?}
}

\cardfoot{Next: outlook and return to Heilbronn, 14:44}

% ---------------------------------------------------
% CARD 9: Outlook, based on Wuttke's notes of 23 November 2026
% ---------------------------------------------------
\cardhead{CARD 9 OF 9 \textbar{} OUTLOOK}{14:44 \textbar{} 1 MIN}{Skilled labour continues to matter}{Capture knowledge, develop it and pass it on}

\cardcolumns{
\lbl{The opportunity, if we get it right}
\begin{cues}
\item Capture the boomers' knowledge
\item Augment it with physical AI
\item Share it with the next generation of skilled workers
\end{cues}

\lbl{What changes}
\handover{New job profiles. New skills to teach and learn.}
}{
\lbl{Return to Heilbronn}
\handover{I am happy to contribute to educating that next generation here at TUM in Heilbronn.}

\lbl{Thank the panel and audience}
\note{Thank you, Helmut, Sarah, Michael and Eike. And thank you all for joining us.}
}

\cardfoot{End: 14:45}

\end{document}
